\chapter{Conclusions and Future Work}
\label{ch:con}
\section{Conclusions}
The research goal of this thesis is development of an end-to-end generalised CRS algorithm that maximises recommendation relevance while conducting a dialogue in a human-like manner and efficiently utilising different modalities of information given by users. Specifically, the focus is on developing an efficient strategy for information elicitation, or obtaining the most valuable information from users given the sparse and unbounded space of possible item features in order to uncover user preferences in the smallest number of conversational turns.

We have proposed a novel CRS system design based on a goal-driven self-supervised bot-play approach. The key components of the system are QBot, ABot, recommendation network $\phi$, and policy network $\pi$. QBot proactively elicits preference information from ABot and updates its internal representation of the ABot's ratings. ABot retrieves the item and attribute preferences from the user's history in response to QBot's questions. The recommendation network predicts movie ratings based on observed preferences, and the policy network guides QBot's behaviour. The collaborative data matrix is used to pre-train the recommendation component, and the change in the recommendation model loss in QBot's internal representation is used as a reward to train the policy model. 

Item attributes, or movie metadata are introduced into all of the individual system components alongside item IDs. This is crucial since users have been found to rely on generalisations and descriptions of movie attributes, such as genre or plot, in addition to sharing examples of movies they liked or disliked. 

We have shown that the recommendation component has successfully learned mappings between items and attributes, which allows the QBot to make accurate rating predictions from a limited set of answered questions. On both synthetic and real world datasets we have shown that trained policy model is more efficient at recreating user preferences than the random baseline. On the synthetic dataset, the policy has learned to minimise recommendation model error in the smallest theoretically possible number of conversational turns. On the MovieLens-based dataset, the policy network has learned to significantly reduce loss compared to the baseline, however, the effect is less pronounced due to weaker item intercorrelations.

The most important implication of work to date is the proposed algorithm that directly addresses the need to develop an efficient active learning strategy for a conversational recommendation system, a problem that does not pertain to classical non-interactive recommenders. To date, this problem is often neglected by authors of similar systems \citep{li2018towards, kang2019recommendation}, which focus on imitating human speech instead. The described approach can be applied on top of supervised pre-training of the encoder-decoder model, which can result in a conversational system that has high practical value in the real world.

The described architecture is comparatively basic, but can be extended to encompass multimodal information obtained from users in a natural language form. This would enable development of a highly flexible system that would be able to ‘understand’ and solicit a wide spectrum of information from users.
Finally, the system is modular and allows to separately pre-train individual components. For example, the full MovieLens dataset can be used to pre-train the policy network, while the encoder-decoder components can be trained on a much more limited natural conversational dataset, such as ReDial \citep{li2018towards}.

The system described in this thesis has a number of limitations. Firstly, the algorithm has only been trained on small datasets to date. In order to scale the approach to large datasets (e.g., full MovieLens dataset) or expand it to natural language inputs, the model architecture needs to be reviewed. The feedforward architecture that is currently used for the policy network does not scale well to large datasets and cannot be used when input space is unbounded (for example, with natural language inputs). Replacing this architecture with a recurrent network with attention mechanism to support arbitrary order of input data would help address this limitation. In addition to that, more advanced models, such as Hierarchical Recurrent Encoder or Memory Network \citep{das2017learning} will be tested for this purpose.

Secondly, the current approach focuses on the accuracy of the representation of user’s tastes, but does not explicitly address the recommendation task.  It could be argued that predicting the full set of user’s ratings is not practical in a recommendation scenario, as it is more important to identify movies that would be rated highly by the user. Redefining the Questioner’s reward to reflect usefulness of the top ranked items (for example, mean reciprocal rank of the top recommendation) would allow to address this issue in the future. Further to that, it is important to consider the novelty aspect of the recommendation -- in a real world scenario it is important to train the agent to discover items that the user has not consumed before.

In conclusion, this thesis proposes a novel end-to-end generalised CRS algorithm that maximises recommendation relevance and is compatible with natural language interfaces. The proposed system design is based on a goal-driven self-supervised bot-play approach. The system has been shown to be efficient in eliciting preference information from users and predicting accurate ratings from a limited set of answered questions. However, the algorithm has only been trained on small datasets and the model architecture needs to be reviewed to scale the approach to larger datasets or expand it to multimodal and natural language inputs. Overall, the proposed algorithm addresses the need to develop an efficient active learning strategy for a conversational recommendation system and has high practical value in the real world.


\section{Future work}

One of the key objectives of any conversational recommendation system is efficient user preference elicitation (acquiring information from a user), as this capability is what distinguishes an interactive system from a traditional recommender engine and allows to build an informative user profile in a cold-start scenario. While some work has been done in this space, this remains a challenging task, to a large extent due to the sparse and unbounded space of features that contribute to user consumption decisions and the lack of rich and natural datasets.

One promising avenue for future research is the application of a self-supervised bot-play approach to conversational recommendation systems. Drawing inspiration from successful techniques developed in other domains, this approach can tackle the problem of user preference elicitation by leveraging bot play. By integrating this approach, conversational recommendation systems can benefit from improved user interaction and information gathering.

Another promising direction for future work involves enhancing dialogue management in conversational recommendation systems by incorporating Language Model-based (LLM-based) models. By extending the system with an LLM-based Dialogue Management module, researchers can improve the system's ability to understand and respond to user inputs effectively. This approach capitalizes on the strengths of LLMs in capturing the intricacies of natural language interactions, enabling more sophisticated and context-aware recommendations. Utilizing pretrained components, such as encoder-decoder models and dialogue management modules trained on massive natural language datasets, is a promising avenue for future work in conversational recommendation systems. By incorporating these pretrained components into the system architecture, researchers can harness the benefits of pretraining and transfer learning, leading to improved performance in recommendation accuracy, user satisfaction, and efficiency.

The bot-play approach for information elicitation module training, coupled with a natural language interface, shows promise as a framework for training conversational recommendation systems that emulate human-like behavior. Future research can delve deeper into the application of this approach, examining its potential in enhancing the quality and relevance of recommendations. By leveraging the bot-playframework and combining it with other techniques, researchers can develop more robust conversational recommendation systems that adapt to users' preferences and needs.